\documentclass[11pt]{article}
% ---------- packages ---------- 
\usepackage[utf8]{inputenc} 
\usepackage[T1]{fontenc} 
\usepackage[margin=1in]{geometry} 
\usepackage{amsmath,amssymb} 
\usepackage{graphicx} 
\usepackage{booktabs} 
\usepackage{array} 
\usepackage{caption} 
\usepackage{xcolor} 
\usepackage{enumitem} 
\usepackage{tikz} 
\usetikzlibrary{positioning,arrows.meta,fit} 
\usepackage[section]{placeins} 
\usepackage[colorlinks=true,linkcolor=blue,urlcolor=blue,citecolor=blue]{hyperref}

\graphicspath{{stage2/results/figures/}{./}} 
\newcommand{\code}[1]{\texttt{#1}}
\usepackage{pifont}
\newcommand{\cmark}{\ding{51}}
\newcommand{\xmark}{\ding{55}}

% ---------- title ---------- 
\title{\textbf{Reproducing, Diagnosing, and Rescuing the CovidGAN Claim}\\[2pt] 
\normalsize Final Project, Part 2 } 
\author{Hila Fishman \and Maya Hayat} 
\date{}

\begin{document} 
\maketitle 
\vspace{-2.0em}

\begin{abstract}
We reconstruct \textbf{CovidGAN} (Waheed et al., 2020), an AC-GAN that synthesizes chest X-rays to augment COVID-19 detection, and test its central claim: a \textbf{+10-point} accuracy lift from synthetic augmentation. Our faithful reimplementation does not reproduce it --- the lift is \textbf{flat ($-0.52$)} against an already-high 90.62\% baseline. Five diagnostic experiments locate two causes: the modern Kaggle data has drifted to become more separable than the paper's 2020 snapshot, and the paper's \emph{frozen} VGG16 backbone leaves only 33K parameters able to respond to new data. Stage 2 attacks both --- a de-fingerprinted generator (noise 1.0, DiffAugment, spectral norm, projection discriminator) and an unfrozen encoder --- and recovers a real lift: \textbf{94.48\% $\pm$ 0.97} to \textbf{96.67\% $\pm$ 0.71} over 5 seeds, past the paper's 95\%. Replicating the pipeline on \textbf{CBIS-DDSM} mammography reproduces every conclusion on independent data, confirming the failure is not an artifact of our COVID set. Our sharpest finding is negative: across four generators, quality metrics (FID, transfer probe) are \emph{inversely} related to downstream value --- the best generator we built produces the only augmentation lift indistinguishable from zero. Capacity, not generator quality, is what governs whether augmentation helps.
\end{abstract}

\section{Original Architecture} \label{sec:original_architecture}
Figure~\ref{fig:arch} gives the block view of the paper's pipeline: an \textbf{AC-GAN} that synthesizes chest X-rays, and a \textbf{VGG16-based classifier} that detects COVID-19 from them. We reconstructed it layer by layer in PyTorch (\code{covidgan/models.py}) and audited every component against the paper's Figures~1--3 and Section~II before interpreting any result. Rather than restate the diagram in prose, we draw out the three specification details that drive everything that follows.

\begin{figure}[htbp]
\centering
\tikzset{
  b/.style={draw,rounded corners,minimum height=7mm,minimum width=14mm,font=\scriptsize,align=center,fill=blue!8},
  io/.style={draw,rounded corners,minimum height=7mm,minimum width=12mm,font=\scriptsize,align=center,fill=black!8},
  frz/.style={draw,rounded corners,minimum height=7mm,minimum width=14mm,font=\scriptsize,align=center,fill=black!18},
}
\textbf{(a) AC-GAN generator}\\[2pt]
\resizebox{\textwidth}{!}{%
\begin{tikzpicture}[>=Latex,node distance=4mm and 6mm]
\node[io] (z) {noise  $z$ \\ $\mathbb{R}^{100}$ };
\node[io,below=of z] (c) {label  $c$ };
\node[b,right=of z] (zd) {Dense+ReLU\\ $7{\times}7{\times}1024$ };
\node[b,right=of c] (cd) {Embed $\to$ Dense\\ $7{\times}7{\times}1$ };
\node[b,right=7mm of zd] (cat) {concat\\ $7{\times}7{\times}1025$ };
\node[b,right=of cat] (u1) {TConv $\uparrow$ \\ $14^2$ };
\node[b,right=of u1] (u2) {TConv $\uparrow$ \\ $28^2$ };
\node[b,right=of u2] (u3) {TConv $\uparrow$ \\ $56^2$ };
\node[b,right=of u3] (u4) {TConv+tanh\\ $112^2$ };
\node[io,right=of u4] (img) {image\\ $112^2{\times}3$ };
\draw[->] (z) -- (zd);
\draw[->] (c) -- (cd);
\draw[->] (zd) -- (cat);
\draw[->] (cd) -- (cat);
\draw[->] (cat) -- (u1);
\draw[->] (u1) -- (u2);
\draw[->] (u2) -- (u3);
\draw[->] (u3) -- (u4);
\draw[->] (u4) -- (img);
\end{tikzpicture}}\\[8pt]
\textbf{(b) AC-GAN discriminator}\hspace{2mm}\textbf{(c) VGG16 detector}\\[2pt]
\resizebox{\textwidth}{!}{%
\begin{tikzpicture}[>=Latex,node distance=4mm and 6mm]
\node[io] (img) {image};
\node[b,right=of img] (d5) {5 conv $\downarrow$ \\ $7{\times}7{\times}512$ };
\node[b,above right=2mm and 7mm of d5,fill=green!10] (val) {validity\\(real/fake)};
\node[b,below right=2mm and 7mm of d5,fill=orange!12] (cls) {class\\(COVID/Normal)};
\draw[->] (img) -- (d5);
\draw[->] (d5) -- (val);
\draw[->] (d5) -- (cls);
\node[io,right=22mm of val] (im2) {image};
\node[frz,right=of im2] (vgg) {VGG16 base\\\textit{frozen}};
\node[b,right=of vgg] (gap) {GAP+Dense(64)\\Dropout};
\node[b,right=of gap] (o2) {Dense(2)\\softmax};
\draw[->] (im2) -- (vgg);
\draw[->] (vgg) -- (gap);
\draw[->] (gap) -- (o2);
\end{tikzpicture}}
\caption{\textbf{(a)}~Generator fuses noise + label and upsamples  $7{\to}112$  through four transpose-conv blocks ($\sim$ 22M params). \textbf{(b)}~Discriminator downsamples to  $7{\times}7{\times}512$  and splits into the two AC-GAN heads ($\sim$ 2M). \textbf{(c)}~Detector is a \emph{frozen} ImageNet VGG16 + small head --- only  $\sim$ 33K of  $\sim$ 14.7M params train, per the paper's original specification.}
\label{fig:arch}
\end{figure}

\begin{itemize}[nosep, leftmargin=1.5em]
  \item \textbf{The latent noise is nearly switched off.} The generator draws $z \in \mathbb{R}^{100}$ from $\mathcal{N}(0, 0.02)$ --- the paper's own value. At $\sigma = 0.02$ the noise vector barely varies, so the output is driven by the \emph{class label} rather than by $z$. This single constant explains both the mode collapse and the class fingerprint we diagnose in Section~\ref{sec:reconstruction_results}, and it is the first thing Track A changes.
  \item \textbf{Conditioning is by label embedding.} A 50-dimensional embedding projected to $7{\times}7{\times}1$ and concatenated with the noise branch's $7{\times}7{\times}1024$ gives the generator an explicit and very cheap channel for encoding class identity.
  \item \textbf{The classifier's backbone is frozen.} Only $\sim$33K of 14.7M parameters train. This is faithful to the paper (Sec.~II-B: ``\ldots setting the `trainable' property on each of the VGG layers to False''), not a shortcut we introduced --- so any anomaly we find is a property of the data or the training regime, not an implementation error.
\end{itemize}

\paragraph{Objective and hyperparameters.} The AC-GAN combines an adversarial term $L_S$ (binary cross-entropy on the validity head, with one-sided label smoothing at 0.9) and an auxiliary classification term $L_C$ (cross-entropy on the class head, applied to real and synthetic alike):
\[ L_S = \mathbb{E}[\log P(S{=}\text{real}|X_{\text{real}})] + \mathbb{E}[\log P(S{=}\text{fake}|X_{\text{fake}})], \qquad
   L_C = \mathbb{E}[\log P(C{=}c|X_{\text{real}})] + \mathbb{E}[\log P(C{=}c|X_{\text{fake}})]. \]
$D$ maximizes $L_S + L_C$ while $G$ maximizes $L_C - L_S$; the detector trains under plain cross-entropy. The paper's hyperparameters: GAN --- batch 64, Adam $\mathrm{lr}=2\times10^{-4}$, $\beta_1 = 0.5$, \textbf{2000 epochs}; classifier --- batch 16, Adam $\mathrm{lr}=10^{-3}$, 15 epochs, images $112\times112$. Data split: train 331 COVID / 601 Normal (932 images), test 72/120 (\textbf{192 images}).

\section{Paper Results} \label{sec:paper_results}
The reported baseline results in Waheed et al. (2020) set a clear target for our reconstruction. 

\subsection{The Baseline Lift Claim}
The paper's main argument is that training a VGG16 classifier on actual chest X-ray data alone yields a relatively low accuracy due to severe data scarcity, which can be dramatically boosted by mixing in synthetic images.
\begin{itemize}[nosep, leftmargin=1.5em]
  \item \textbf{CNN-AD (Actual Data Only):} 85.00\% accuracy.
  \item \textbf{CNN-SA (Synthetic Augmented):} 95.00\% accuracy (a massive \textbf{+10.00\%} absolute lift).
  \item \textbf{Recall on COVID Class:} The paper's CNN-AD baseline suffered from a low COVID-class recall of 0.69, missing 22 out of 72 real COVID cases. Augmentation raised this recall to 0.90, which represents a clinically vital improvement.
\end{itemize}

% \subsection{Metrics and What They Measure}
% \begin{itemize}[nosep, leftmargin=1.5em]
%   \item \textbf{Accuracy} \emph{(the paper's metric)} --- correct/total on the same 192-image real test set, the number the +10-point claim is stated in.
%   \item \textbf{Per-class recall} \emph{(studied in class)} --- $\mathrm{TP}/(\mathrm{TP}{+}\mathrm{FN})$. Missing a COVID case costs more than a false alarm, so recall is the clinically meaningful axis --- and it is where the paper locates its gain.
%   \item \textbf{FID} \emph{(studied in class, for the generator)} --- Fr\'{e}chet distance between InceptionV3 features of real vs.\ synthetic images; lower is better. Biased upward on a test set as small as ours, so we read it relatively across our own generators.
%   \item \textbf{Transfer probe} \emph{(ours)} --- train on synthetic \emph{alone}, test on real. Reference line is the \textbf{62.5\% majority-class floor}; clearing it means the fakes carry transferable signal, landing below it means they cut against it.
% \end{itemize}

\noindent A theme from Section~\ref{sec:reconstruction_results} onward: the generator-side measures and the downstream lift can move in opposite directions, so only a direct CNN-AD vs.\ CNN-SA comparison tests the paper's claim.

\subsection{Metrics Used}
\begin{itemize}[nosep, leftmargin=1.5em]
  \item \textbf{FID} \emph{(studied in class, for the generator)} --- Fr\'{e}chet distance between InceptionV3 features of real vs.\ synthetic images; lower is better. Biased upward on a test set as small as ours, so we read it relatively across our own generators.
  \item \textbf{Transfer probe} \emph{(ours)} --- train on synthetic \emph{alone}, test on real. Reference line is the \textbf{62.5\% majority-class floor}; clearing it means the fakes carry transferable signal, landing below it means they cut against it.
\end{itemize}


\section{Reconstruction Results (Stage 1)} \label{sec:reconstruction_results}
When we created and executed a layer-by-layer reconstruction of the entire pipeline, we obtained results that did not match the paper's claims. 

\subsection{The Reconstruction Anomaly}
Two main anomalies stood out when running the models: (A) our baseline classifier (CNN-AD) is already much stronger than the paper's baseline (90.62\% vs. 85.00\%), and (B) adding synthetic images to our classifier yields a completely flat result (-0.52\% lift), meaning the augmentation effect vanishes. To isolate the root causes, we designed a series of five controlled diagnostic experiments.

\subsection{The Five Controlled Diagnostic Experiments}

\paragraph{Test 1: Implementation audit.} A parameter-count and layer-freeze audit against the paper's specification matched exactly, ruling out an implementation error.

\paragraph{Test 2: Cross-source shortcut }
Our reconstruction draws COVID images from the IEEE database and Normal images from Kaggle exactly as the paper did. But mixing sources means a classifier could separate the classes on processing cues rather than pathology. To test this we rebuilt the split with both classes drawn from the single Kaggle Radiography Database. Accuracy remained unchanged, meaning that forcing a single source did not pull the baseline toward 85\%.

\paragraph{Test 3: DDSM Dataset} A central concern in our study is the data specificity, in other words, is our failure to reproduce the paper's accuracy increase an artifact of the modern, cleaner COVID-CXR dataset, or is it a property of the AC-GAN method itself under frozen-backbone constraints? To answer this, we
replicated the entire pipeline on \textbf{CBIS-DDSM}, which is a well known breast mammography dataset containing both mass and calcification lesions. In this experiment we focused on the mass ROI image lesions, containing 1,318 real training images and 378 test images (231 benign, 147 malignant).

\paragraph{Data adjustment.} DDSM ROIs are stored as DICOM and converted to 16-bit grayscale PNGs with
pixel values up to 65,535. Standard PIL conversion truncates to 8 bits, collapsing nearly every pixel
to flat white. We implemented per-image min--max scaling to $[0,255]$ prior to RGB conversion,
restoring tissue texture and rescuing the baseline classifier. All results below post-date this fix.

\begin{table}[htbp]
\centering
\caption{CBIS-DDSM main results. Each row averages three independent runs. Balanced accuracy is the
mean of malignant and benign recall.}
\label{tab:ddsm_main}
\begin{tabular}{llrrrrr}
\toprule
\textbf{Mode} & \textbf{Classifier} & \textbf{Train acc} & \textbf{Test acc} &
\textbf{Mal.\ recall} & \textbf{Ben.\ recall} & \textbf{Bal.\ acc} \\
\midrule
\textbf{AD (real only)} & frozen & 0.795 & \textbf{66.05} & 0.535 & 0.740 & \textbf{0.638} \\
SA (baseline GAN) & frozen & 0.848 & 63.80 & 0.556 & 0.691 & 0.623 \\

\bottomrule
\end{tabular}
\end{table}

\paragraph{Test 4: Coarse-resolution shortcut}
Even within one source, the model might be reading coarse global properties (brightness, framing, silhouette) rather than pathology. We downsampled real images to $N \times N$ and upsampled back to $112 \times 112$, destroying fine anatomical detail while preserving coarse structure, and retrained CNN-AD at each scale (Table~\ref{tab:coarse_resolution}). A purely coarse shortcut would survive this, holding near 89\% even at $8 \times 8$. Instead accuracy falls 16 points and COVID recall roughly halves --- the collapse is specific to COVID. This weighs against a coarse shortcut and points to a model reading genuine fine detail, with a small non-anatomical residual remaining (at $4 \times 4$, where no anatomy survives, accuracy is still 11 points above the 62.5\% floor).

\begin{table}[htbp]
\centering
\caption{\textbf{Coarse-resolution shortcut check.} Downsampling real images to $N \times N$ and upsampling back to 112 destroys fine anatomical details but retains coarse global properties. Classifier accuracy decays, proving detail-dependence.}
\label{tab:coarse_resolution}
\begin{tabular}{@{}ccc@{}}
\toprule
\textbf{Resolution ($N \times N$)} & \textbf{Accuracy} & \textbf{COVID Recall} \\
\midrule
112 (Full) & 89.58\% & 0.83 \\
32 & 83.85\% & 0.64 \\
16 & 82.29\% & 0.61 \\
8 & 77.08\% & 0.51 \\
4 & 73.96\% & 0.51 \\
\bottomrule
\end{tabular}
\end{table}

\paragraph{Test 5: Synthetic-only probe (do the fakes carry pathology?)}
This test began as an observation rather than a hypothesis: looking at the generated samples, the difference between the classes is plainly visible to the eye. In Figure~\ref{fig:fingerprint} the synthetic COVID tiles are uniformly brighter and hazier than the darker, sharper Normal tiles --- a whole-image style difference, which is emphatically not how real pathology presents. That prompted the question of whether the synthetic class signal is pathology at all, which the probe answers numerically: we train the classifier on the synthetic pool alone and test on real images.

\begin{figure}[htbp]
\centering
\includegraphics[width=0.72\linewidth]{gan_samples.jpg}
\caption{\textbf{The fingerprint, seen.} 2000-epoch AC-GAN samples: COVID tiles (columns 1 and 3) are uniformly brighter and hazier than the darker, sharper Normal tiles, so the classes separate on \emph{global brightness style} rather than localized pathology.}
\label{fig:fingerprint}
\end{figure}

It reaches nearly 100\% on synthetic training data but only \textbf{65.45\% $\pm$ 0.80} on the real test set (3 seeds) --- barely three points above the 62.5\% majority-class floor. Near-perfect separability inside the synthetic set therefore buys almost nothing on real images: the AC-GAN stamps an easy, class-conditioned signature through its label branch, trivially recoverable within its own output but only weakly related to what separates real COVID from real Normal. An earlier 25-epoch ``smoke'' generator makes the point more starkly still --- its pool probed at \textbf{55.21\%} (COVID recall 0.22), \emph{below} the floor, meaning a model trained only on those images did worse than guessing all-Normal.

\paragraph{Test 6: Generator Quality Over Training (FID)}
We computed the overall and per-class FID scores between our real test set and the synthetic pool at an early epoch (epoch 25) vs. the final epoch (epoch 2000), shown in Table~\ref{tab:fid}. The FID roughly halved over training, indicating the generator did learn to produce CXR-like structures (visible ribcages and lung fields). Yet mixing these markedly better images into the classifier yielded an essentially identical flat result (90.10\% vs. 90.62\%). This is our strongest evidence on Anomaly B: if image \emph{quality} were the binding constraint, a $2\times$ improvement in FID should have produced some downstream movement, and it produced almost none. We therefore treat image quality as unlikely to be what caps augmentation here --- strong support rather than proof, since FID is an imperfect proxy on so few real images.


\begin{table}[htbp]
\centering
\caption{\textbf{Fr\'{e}chet Inception Distance (FID) Over Training.} FID scores computed between the real test set and the synthetic pool at early vs. final epoch, indicating that generator image quality is not what is capping augmentation.}
\label{tab:fid}
\begin{tabular}{@{}lccc@{}}
\toprule
\textbf{Set} & \textbf{25-Epoch Generator} & \textbf{2000-Epoch Generator} & \textbf{Net Change} \\
\midrule
Overall & 504.4 & \textbf{272.7} & \textbf{-46.0\%} \\
COVID & 487.2 & \textbf{302.2} & \textbf{-38.0\%} \\
Normal & 538.0 & \textbf{290.2} & \textbf{-46.0\%} \\
\bottomrule
\end{tabular}
\end{table}

\subsection{The Visual Fingerprint and Anomaly Explanations}
Reading the five tests together, a coherent explanation for both anomalies emerges:
\begin{itemize}[nosep, leftmargin=1.5em]
  \item \textbf{Why our baseline is higher (Anomaly A):} The modern Kaggle radiography dataset, while matching the original sources, has been repeatedly cleaned, updated, and enlarged since the paper's 2020 release. Our classifier is performing legitimate, detail-based classification on a cleaner, much more separable dataset, saturating its performance and leaving no headroom for a +10\% lift.
  \item \textbf{Why augmentation is flat (Anomaly B):}
  \begin{enumerate}[label=\arabic*., nosep, leftmargin=1.5em]
    \item \textbf{The Visual Fingerprint:} When we visually inspected the generated samples (Figure~\ref{fig:fingerprint}), the reason for the flat results became immediately apparent. The synthetic COVID chest X-rays were \textbf{conspicuously brighter, hazier, and lower-contrast} than the Normal fakes, which were uniformly dark and sharp.
    We ran a script to quantify this: on \textit{real} data, the class-mean brightness difference is a tiny \textbf{+8.8 out of 255 (only $\sim$3.5\%)}, with overlapping ranges. The AC-GAN, finding the path of least resistance to satisfy its classification loss, amplified this tiny, overlapping real tendency into a \textbf{deterministic global brightness shortcut}. It stamped a global brightness signature instead of learning localized lung pathology.
    \item \textbf{The Frozen Classifier Bottleneck:} Because the VGG16 backbone is completely frozen, the classifier only trains its $\sim$33K head parameters. It cannot adapt or learn new features from the synthetic data; it can only nudge a linear decision boundary on top of pre-fixed ImageNet representations. This frozen head acts as a ceiling on how much augmentation can contribute --- and, symmetrically, on how much the weak synthetic signal can \emph{hurt}, which is consistent with the observed result being flat rather than negative.
  \end{enumerate}
\end{itemize}

\paragraph{An open caveat.} Test 3 argues against a \emph{coarse} shortcut, but it cannot rule out a \textbf{fine-scale, dataset-specific artifact} --- a rescaling-kernel signature, a JPEG or compression fingerprint, or a faint watermark surviving from the source archive. Such an artifact would live in exactly the fine detail that downsampling destroys, so it would produce the same decay pattern in Table~\ref{tab:coarse_resolution} that genuine pathology does, and our blurring test could not tell the two apart. We therefore state Anomaly A's explanation as the reading most consistent with the evidence, not as a settled result. Settling it would require evaluating CNN-AD on an independent public CXR dataset carrying both classes and sharing no images with our sources; we discuss why that is harder than it sounds in Section~\ref{sec:discussion}.

\section{Improved Architecture (Stage 2)} \label{sec:improved_architecture}
Our Stage 2 architecture is a \textbf{coupled two-front redesign} that unblocks the generator and the classifier in turn. First, we built a better generator and after finding that it still did not increase the accuracy, we decided to focus on the classifier (Track B) and discover that the small, frozen model was the real bottleneck all along. That sequence is itself one of the findings of this project.

\subsection{Track A: Generator Improvements (Killing the Fingerprint)}
To prevent the generator from collapsing into a deterministic ``brightness fingerprint'' and to stabilize training, we made four changes:
\begin{enumerate}[nosep, leftmargin=1.5em]
  \item \textbf{Unlocking Latent Noise ($z \sim \mathcal{N}(0, 1)$):} The original paper's code used a noise standard deviation of $0.02$, which practically turned the noise vector off, forcing the generator to collapse to a single prototype per class. Raising this to $1.0$ restored sample diversity.
  \item \textbf{Spectral Normalization:} We wrapped every convolutional layer in the discriminator and both of its output heads with spectral normalization to bound the Lipschitz constant, preventing discriminator explosion on small datasets.
  \item \textbf{DiffAugment:} We applied differentiable translation, color, and cutout augmentations to both real and fake images during GAN training, stabilizing the discriminator gradients.
  \item \textbf{Projection Discriminator:} To kill the class fingerprint at its root, we completely replaced the AC-GAN's auxiliary classifier head. The \textbf{Projection Discriminator} (Miyato et al., 2018) conditions the discriminator via an inner product between a class embedding and the intermediate feature representation, removing the incentive to embed a cheap, global class marker.
\end{enumerate}

We trained two variants: an \textbf{Improved AC-GAN} (using changes 1-3) and a \textbf{Projection GAN} (adding change 4), both for a highly efficient 300 epochs.

% =====================================================================
% Figure: Stage 2 improved GAN architecture
% Requires (already in the preamble): \usepackage{tikz}
%                                     \usetikzlibrary{positioning,arrows.meta,fit}
% =====================================================================
\begin{figure}[htbp]
\centering
% ---- shared node styles (same palette as Fig.~\ref{fig:arch}) ----
\tikzset{
  b/.style={draw,rounded corners,minimum height=8mm,minimum width=15mm,font=\scriptsize,align=center,fill=blue!8},
  io/.style={draw,rounded corners,minimum height=8mm,minimum width=13mm,font=\scriptsize,align=center,fill=black!8},
  new/.style={draw,very thick,rounded corners,minimum height=8mm,minimum width=15mm,font=\scriptsize,align=center,fill=red!12},
  op/.style={draw,circle,inner sep=1pt,font=\scriptsize,fill=black!5},
}

% ---- (a) Improved generator ----
\textbf{(a) Improved generator --- architecture unchanged, noise scale fixed}\\[2pt]
\resizebox{\textwidth}{!}{%
\begin{tikzpicture}[>=Latex,node distance=5mm and 7mm]
  \node[new] (z) {noise $z$\\$\mathcal{N}(0,\mathbf{1})$};
  \node[io,below=of z] (c) {label $c$};
  \node[b,right=of z] (zd) {Dense+ReLU\\$7{\times}7{\times}1024$};
  \node[b,right=of c] (cd) {Embed$\to$Dense\\$7{\times}7{\times}1$};
  \node[b,right=8mm of zd] (cat) {concat\\$7{\times}7{\times}1025$};
  \node[b,right=of cat] (u1) {TConv$\uparrow$\\$14^2$};
  \node[b,right=of u1] (u2) {TConv$\uparrow$\\$28^2$};
  \node[b,right=of u2] (u3) {TConv$\uparrow$\\$56^2$};
  \node[b,right=of u3] (u4) {TConv+tanh\\$112^2$};
  \node[io,right=of u4] (img) {image\\$112{\times}112{\times}3$};
  \draw[->] (z) -- (zd); \draw[->] (c) -- (cd);
  \draw[->] (zd) -- (cat); \draw[->] (cd) -- (cat);
  \draw[->] (cat) -- (u1); \draw[->] (u1) -- (u2); \draw[->] (u2) -- (u3);
  \draw[->] (u3) -- (u4); \draw[->] (u4) -- (img);
\end{tikzpicture}}\\[10pt]

% ---- (b) Improved AC-GAN discriminator ----
\textbf{(b) Improved AC-GAN discriminator --- best \emph{downstream} result}\\[2pt]
\resizebox{\textwidth}{!}{%
\begin{tikzpicture}[>=Latex,node distance=5mm and 7mm]
  \node[io] (img) {real / fake\\image $112^2{\times}3$};
  \node[new,right=of img] (aug) {DiffAugment\\{\tiny color+translation+cutout}};
  \node[new,right=of aug] (d1) {SN-Conv$\downarrow$};
  \node[new,right=of d1] (d2) {SN-Conv$\downarrow$};
  \node[new,right=of d2] (d3) {SN-Conv$\downarrow$};
  \node[new,right=of d3] (d4) {SN-Conv$\downarrow$};
  \node[new,right=of d4] (d5) {SN-Conv$\downarrow$\\$7{\times}7{\times}512$};
  \node[b,right=of d5] (fl) {Flatten\\$25088$};
  \node[b,above right=3mm and 9mm of fl,fill=green!10] (val) {validity head\\(real / fake)};
  \node[b,below right=3mm and 9mm of fl,fill=orange!12] (cls) {class head\\(COVID / Normal)};
  \draw[->] (img) -- (aug); \draw[->] (aug) -- (d1); \draw[->] (d1) -- (d2);
  \draw[->] (d2) -- (d3); \draw[->] (d3) -- (d4); \draw[->] (d4) -- (d5);
  \draw[->] (d5) -- (fl);
  \draw[->] (fl) -- (val); \draw[->] (fl) -- (cls);
\end{tikzpicture}}\\[10pt]

% ---- (c) Projection discriminator ----
\textbf{(c) Projection discriminator --- best \emph{generator} (no auxiliary classifier)}\\[2pt]
\resizebox{\textwidth}{!}{%
\begin{tikzpicture}[>=Latex,node distance=5mm and 7mm]
  \node[io] (img) {real / fake\\image $112^2{\times}3$};
  \node[new,right=of img] (aug) {DiffAugment\\{\tiny color+translation+cutout}};
  \node[new,right=of aug] (d1) {SN-Conv$\downarrow$};
  \node[new,right=of d1] (d2) {SN-Conv$\downarrow$};
  \node[new,right=of d2] (d3) {SN-Conv$\downarrow$};
  \node[new,right=of d3] (d4) {SN-Conv$\downarrow$};
  \node[new,right=of d4] (d5) {SN-Conv$\downarrow$\\$7{\times}7{\times}512$};
  \node[b,right=of d5] (pool) {ReLU +\\global sum pool\\$h\in\mathbb{R}^{512}$};
  \node[b,above right=4mm and 10mm of pool,fill=green!10] (val) {$w^{\!\top}h+b$\\(real / fake)};
  \node[new,below right=4mm and 10mm of pool] (ip) {$\langle \mathbf{e}(y),\,h\rangle$\\projection};
  \node[io,below=14mm of d3] (y) {label $y$};
  \node[b,right=16mm of y,fill=orange!12] (emb) {SN-Embed\\$2{\times}512$};
  \node[op,right=32mm of pool] (sum) {$+$};
  \node[io,right=of sum] (out) {$D(x,y)$\\{\tiny hinge loss}};
  \draw[->] (img) -- (aug); \draw[->] (aug) -- (d1); \draw[->] (d1) -- (d2);
  \draw[->] (d2) -- (d3); \draw[->] (d3) -- (d4); \draw[->] (d4) -- (d5);
  \draw[->] (d5) -- (pool);
  \draw[->] (pool) -- (val); \draw[->] (pool) -- (ip);
  \draw[->] (y) -- (emb); \draw[->] (emb) -- (ip);
  \draw[->] (val) -- (sum); \draw[->] (ip) -- (sum);
  \draw[->] (sum) -- (out);
\end{tikzpicture}}
\caption{Block view of the Stage 2 improved GAN; \textbf{red outlines mark every
change from the Stage 1 reconstruction} (Fig.~\ref{fig:arch}), all of them in the
\emph{training recipe} rather than the layer shapes --- the generator keeps its
$\sim$22M parameters and the discriminator trunk its $\sim$2M, so any gain is
attributable to the recipe, not to a bigger model. \textbf{(a)}~The generator is
architecturally identical to Stage 1; the single change is the latent scale,
$z\sim\mathcal{N}(0,0.02)\to\mathcal{N}(0,1)$, which restores sample diversity and
dilutes the label fingerprint. \textbf{(b)}~The improved AC-GAN discriminator adds
DiffAugment on both real and fake inputs (small-data regularisation) and spectral
normalisation on every conv, dropping Stage 1's Dropout; the twin validity/class
heads are retained for a controlled comparison. \textbf{(c)}~The projection
discriminator shares that trunk exactly but replaces the auxiliary class head with
a projection inner product, $D(x,y)=w^{\!\top}h+b+\langle \mathbf{e}(y),h\rangle$
(Miyato \& Koyama, 2018), trained with the hinge loss --- with no classifier head
to game, the generator can raise $D$ only by matching class $y$'s real features.
The two heads in (b) and (c) are the \emph{only} difference between the two Stage 2
GANs. (c) is the better \emph{generator} (FID $121.3$, transfer probe $81.3\%$ at
600 epochs, vs.\ $225.7$ / $72.9\%$ for (b)), yet (b) gives the better
\emph{downstream} accuracy ($97.92\%$ vs.\ $97.40\%$ with the unfrozen classifier)
--- the quality/utility paradox of \S\ref{sec:trackBresults}.}
\label{fig:arch-improved}
\end{figure}

\subsection{Track B: Classifier Improvements (Unfreezing the Backbone)}
Once we developed a de-fingerprinted generator, we addressed the classifier bottleneck:
\begin{enumerate}[nosep, leftmargin=1.5em]
  \item \textbf{Unfreezing Top Convolutional Blocks:} We unfroze the top blocks of the pre-trained ImageNet VGG16 backbone to allow domain-specific fine-tuning. We evaluated two levels of capacity:
  \begin{itemize}[nosep, leftmargin=1.5em]
    \item Unfreeze Blocks = 1 (uf=1): Unlocks VGG16's last block, releasing 7.11 million trainable parameters.
    \item Unfreeze Blocks = 2 (uf=2): Unlocks the last two blocks, releasing 13.01 million trainable parameters.
  \end{itemize}
  \item \textbf{Adding a BatchNorm Head:} We modified the custom head to stabilize training: \code{Dense(64) -> BatchNorm1d -> ReLU -> Dropout(0.5) -> Dense(2)}.
  \item \textbf{Discriminative Learning Rates:} To prevent destroying pre-trained ImageNet features, we trained the new head at a learning rate of $1 \times 10^{-3}$, while training the unfrozen, highly sensitive VGG16 backbone blocks at a gentle $1 \times 10^{-5}$.
\end{enumerate}

\section{Improved Results (Stage 2)} \label{sec:improved_results}
By combining these improvements, we conducted a series of  experiments that successfully restored the core claims of the paper.

\subsection{Main Result: Paper vs. Reconstruction vs. Improved}
Our final comparative results are presented in Table~\ref{tab:main_results} and visualized in Figure~\ref{fig:comparison}. For Stage 2, we report the mean $\pm$ standard deviation over 5 independent seeds to ensure statistical rigor. Note that this table isolates the \emph{classifier} improvement: every row uses the same Stage-1 synthetic pool, so the only thing changing is encoder capacity. The effect of changing the \emph{generator} is reported separately in Table~\ref{tab:generator_performance}, and the interaction between the two in Table~\ref{tab:reversal_paradox}.

\begin{figure}[htbp]
\centering
\includegraphics[width=0.80\linewidth]{comparison_bar.png}
\caption{\textbf{Paper vs.\ reconstruction vs.\ improved.} Stage 1 (frozen) shows no augmentation gain; both Stage 2 variants raise the baseline and restore a positive lift, with $uf{=}2$ CNN-SA passing the paper's 95\%.}
\label{fig:comparison}
\end{figure}

\begin{table}[htbp]
\centering
\caption{\textbf{The full capacity $\times$ generator grid.} Every pool under the paper's frozen detector (3 seeds) and our unfrozen $uf{=}2$ encoder (5 seeds); lifts paired by seed against that column's own CNN-AD. Under the frozen head no pool achieves a positive lift and two significantly harm the classifier; under the unfrozen encoder every pool helps.}
\label{tab:reversal_paradox}
\resizebox{\linewidth}{!}{%
\begin{tabular}{@{}lcccc@{}}
\toprule
& \multicolumn{2}{c}{\textbf{Frozen head (paper's, 3 seeds)}} & \multicolumn{2}{c}{\textbf{Unfrozen $uf{=}2$ (5 seeds)}} \\
\cmidrule(lr){2-3}\cmidrule(lr){4-5}
\textbf{Synthetic pool} & \textbf{Accuracy} & \textbf{Lift} & \textbf{Accuracy} & \textbf{Lift} \\
\midrule
--- (real only, CNN-AD) & 91.15\% $\pm$ 0.43 & --- & 94.48\% $\pm$ 0.97 & --- \\
\midrule
Baseline AC-GAN, 2000 ep & \textbf{91.32\% $\pm$ 0.30} & $+0.17$ \emph{(n.s.)} & 96.67\% $\pm$ 0.71 & $+2.19$ \\
Baseline AC-GAN, 300 ep & 89.76\% $\pm$ 0.30 & $\mathbf{-1.39}$ & 96.25\% $\pm$ 1.01 & $+1.77$ \\
Improved AC-GAN, 300 ep & 90.28\% $\pm$ 0.30 & $-0.87$ & \textbf{96.98\% $\pm$ 0.61} & $\mathbf{+2.50}$ \\
Projection, 300 ep & 90.97\% $\pm$ 1.59 & $-0.17$ \emph{(n.s.)} & 96.04\% $\pm$ 0.85 & $+1.56$ \\
Projection, 600 ep & 89.76\% $\pm$ 0.30 & $\mathbf{-1.39}$ & 94.79\% $\pm$ 1.09 & $+0.31$ \emph{(n.s.)} \\
\bottomrule
\end{tabular}%
}
\end{table}

\begin{table}[htbp]
\centering
\caption{\textbf{Generator quality versus downstream value --- the reversal.} Left block scores each generator on its own (FID lower is better; transfer probe over 3 seeds, floor $=62.5\%$). Right block feeds that same pool to an identical $uf{=}2$ classifier over 5 seeds, against a common CNN-AD of 94.48\% $\pm$ 0.97. Down the four trained generators every generator-side measure improves while the downstream lift decays to nothing. The 25-epoch smoke row is a reference point only, never used for augmentation.}
\label{tab:generator_performance}
\resizebox{\linewidth}{!}{%
\begin{tabular}{@{}lccccc@{}}
\toprule
& & \multicolumn{2}{c}{\textbf{Generator-side quality}} & \multicolumn{2}{c}{\textbf{Downstream value ($uf{=}2$, 5 seeds)}} \\
\cmidrule(lr){3-4}\cmidrule(lr){5-6}
\textbf{Generator Variant} & \textbf{Epochs} & \textbf{FID} & \textbf{Transfer Probe} & \textbf{CNN-SA} & \textbf{Aug. Lift} \\
\midrule
Smoke AC-GAN (25 ep, \emph{reference}) & 25 & 504.4 & 55.21\% \emph{(below floor)} & --- & --- \\
\midrule
Stage 1 AC-GAN (Baseline) & 2000 & 272.7 & 65.45\% $\pm$ 0.80 & 96.67\% $\pm$ 0.71 & \textbf{+2.19} \\
Stage 1 AC-GAN (Baseline) & \emph{300} & --- & --- & 96.25\% $\pm$ 1.01 & +1.77 \\
Improved AC-GAN & 300 & 225.7 & 72.40\% $\pm$ 0.52 & \textbf{96.98\% $\pm$ 0.61} & \textbf{+2.50} \\
Projection Discriminator & 300 & 155.0 & 74.65\% $\pm$ 0.79 & 96.04\% $\pm$ 0.85 & +1.56 \\
Projection Discriminator & 600 & \textbf{121.3} & \textbf{82.29\% $\pm$ 1.04} & 94.79\% $\pm$ 1.09 & +0.31 \emph{(n.s.)} \\
\bottomrule
\end{tabular}%
}
\end{table}

\begin{table}[htbp]
\centering
\caption{\textbf{Main Results: Paper, Stage 1 Reconstruction, and Stage 2 Improvements.} Accuracy results are evaluated on the same 192-image real test set. Stage 2 improved configurations report the mean and standard deviation over 5 independent seeds.}
\label{tab:main_results}
\resizebox{\linewidth}{!}{%
\begin{tabular}{@{}lcccc@{}}
\toprule
\textbf{Model / Configuration} & \textbf{Trainable Params} & \textbf{CNN-AD (Real Only)} & \textbf{CNN-SA (Augmented)} & \textbf{Augmentation Lift} \\
\midrule
Original Paper & $\sim$33K & 85.00\% & 95.00\% & \textbf{+10.00\%} \\
Stage 1 Reconstruction & $\sim$33K & 90.62\% & 90.10\% & \textbf{-0.52\%} (Flat) \\
Stage 2 Improved ($uf=1$) & $\sim$7.11M & 92.50\% $\pm$ 0.71\% & 93.65\% $\pm$ 1.16\% & \textbf{+1.15\%} \\
Stage 2 Improved ($uf=2$) & \textbf{$\sim$13.01M} & \textbf{94.48\% $\pm$ 0.97\%} & \textbf{96.67\% $\pm$ 0.71\%} & \textbf{+2.19\%} \\
\bottomrule
\end{tabular}%
}
\end{table}

Unfreezing the encoder successfully adapted the ImageNet features to CXR textures, raising our baseline to \textbf{94.48\%}. More importantly, unfreezing capacity \textbf{unlocked the augmentation lift}, allowing the synthetic images to shape the feature representations and achieving an outstanding augmented accuracy of \textbf{96.67\%—comfortably exceeding the paper's 95\% limit}.

\paragraph{COVID-Class Recall Focus:}
We also tracked our clinically vital COVID-class recall across the configurations:
\begin{itemize}[nosep, leftmargin=1.5em]
  \item Stage 2 (uf=1): AD Recall = 88.89\% $\pm$ 1.76\% $\rightarrow$ SA Recall = \textbf{91.39\% $\pm$ 1.62\%} (Lift: \textbf{+2.50\%})
  \item Stage 2 (uf=2): AD Recall = 93.33\% $\pm$ 3.22\% $\rightarrow$ SA Recall = \textbf{96.67\% $\pm$ 2.08\%} (Lift: \textbf{+3.33\%})
\end{itemize}

\subsection{Track A: Generator Performance and the Reversal Paradox}
Our Track A improvements produced large gains on every generator-side measure, as shown in the left half of Table~\ref{tab:generator_performance}. The Projection Discriminator at 600 epochs reaches an \textbf{FID of 121.3} and a \textbf{transfer-probe accuracy of 82.29\% $\pm$ 1.04} --- less than half the Stage-1 FID, and a probe result that clears the 62.5\% floor by nearly 20 points, against the Stage-1 generator's barely-three. By every metric available for judging a generator on its own, the fingerprint is gone and the synthetic images now carry genuine, transferable structure.

\paragraph{The Reversal Paradox (Our Central Finding):}
We then fed each of these four synthetic pools to the \emph{same} $uf{=}2$ classifier, changing nothing but the pool, and ran every arm over five seeds. The right half of Table~\ref{tab:generator_performance} reports the result, and it runs precisely opposite to expectation: \textbf{as the generator improves, the downstream benefit decays.} The weakest generator of the four (Stage-1 AC-GAN: FID 272.7, probe 65.5\%, barely clearing the floor) delivers a $+2.19$-point lift, while the strongest (Projection-600: FID 121.3, probe 82.3\%) delivers $+0.31$ --- the only pool of the four whose lift is statistically indistinguishable from zero.

We want to be careful about what this does and does not establish. The pairwise step from Projection-300 to Projection-600 is $-1.25$ points, which against a paired seed-to-seed spread of $2.07$ ($t \approx 1.35$, $n{=}5$) we cannot call significant on its own; an earlier single-seed run had suggested a much larger $-3.65$, and multi-seeding cut it down. What \emph{is} solid is the trend across all four pools, and the contrast between the two ends of it: the improved AC-GAN outperforms Projection-600 downstream by $+2.19$ points ($t \approx 3.63$, paired, $n{=}5$) despite being the clearly worse generator by FID and by transfer probe. Read together, the four rows make the point far more convincingly than any single pair could: \textbf{generator-side quality metrics do not predict downstream augmentation value, and can invert it.} Our leading explanation is a dilution or feature-space interaction effect --- a more faithful generator produces images that overlap the real distribution more closely, and so displace real examples rather than complementing them --- but we flag this as a hypothesis we have not yet isolated.

\paragraph{A matched-epoch control.} Because the Stage-1 generator trained for 2000 epochs and the improved ones for 300, budget was confounded with architecture. Retraining the \emph{baseline} at exactly 300 epochs (second row of Table~\ref{tab:generator_performance}) removes it: at matched budget the baseline reaches $+1.77$, so neither improved generator significantly separates from it ($+0.73$, $t \approx 1.20$ for the improved AC-GAN; $-0.21$, $t \approx 0.24$ for the projection discriminator). Training the baseline from 300 to 2000 epochs is worth only $+0.42$ ($t \approx 0.52$). Every Track A intervention did what it was designed to do \emph{to the generator}; none moved the downstream number.

\paragraph{Both fixes are necessary --- and only one is doing real work.} Table~\ref{tab:reversal_paradox} crosses every pool against both classifier regimes. Under the frozen head \textbf{no pool achieves a positive lift}; the best manages $+0.17$, and two significantly degrade the classifier ($-1.39$ each). Under the unfrozen encoder the identical pools all help. Read by rows, moving frozen $\to$ unfrozen is worth 5--6 points on every row, while moving between generators is worth at most 2. \textbf{Capacity is the precondition, and on this dataset it is also the dominant term} --- the paper credits the generator; our evidence credits the classifier's ability to respond at all.

% \paragraph{Where CovidGAN and DDSM diverge.} This is the one place our two datasets disagree. On DDSM the same class of intervention produced a large matched-epoch gain (+6.9 points, Table~\ref{tab:ddsm_matched_epoch}); on COVID it produces essentially none. This is consistent with the headroom argument --- DDSM's classifier sits near 64\%, far from ceiling, while ours is already at 94.5\% --- but it does mean neither dataset supports ``a better GAN yields a better classifier'' unconditionally.

\subsection{Supporting Analysis A: Data-Scarcity Curve}
Why is our full-data lift $+2.19$ rather than the paper's $+10$? We subsampled the real training set to 10/25/50\% while holding the synthetic pool fixed ($uf{=}2$, 3 seeds; Table~\ref{tab:scarcity_results}).

\begin{table}[htbp]
\centering
\caption{\textbf{Data-Scarcity Sweep Results ($uf=2$).} Classification accuracy across different real training set fractions with a fixed synthetic pool (evaluated over 3 seeds).}
\label{tab:scarcity_results}
\resizebox{\linewidth}{!}{%
\begin{tabular}{@{}ccccc@{}}
\toprule
\textbf{Real Data Fraction} & \textbf{Real Train Size ($N_{\text{real}}$)} & \textbf{CNN-AD (Real Only)} & \textbf{CNN-SA (Augmented)} & \textbf{Absolute Accuracy Lift} \\
\midrule
10\% & 93 & 86.98\% & 89.24\% & \textbf{+2.26\%} \\
25\% & 233 & 89.58\% & 92.88\% & \textbf{+3.30\%} (Peak Lift) \\
50\% & 466 & 91.32\% & 93.06\% & \textbf{+1.74\%} \\
100\% & 932 & 93.92\% & 96.01\% & \textbf{+2.08\%} \\
\bottomrule
\end{tabular}%
}
\end{table}

The baseline degrades monotonically as data shrinks (93.9 $\to$ 87.0\%) while augmentation stays positive at every level, peaking at $+3.30$ around 25\%. This directionally supports the paper's ``augmentation rescues data-starved models'' thesis and explains our smaller full-data lift as a ceiling effect. The curve is not perfectly monotonic and the 10\% point is noisy (93 images; its seeds ran $+4.17$, $-1.04$, $+3.65$), so we call it consistently positive and strongest in the low-to-mid regime rather than monotonic.

\paragraph{The same sweep under a frozen head reverses sign.} Repeating it with the paper's frozen detector and our best (Projection) pool turns every gain into a loss, worsening as real data shrinks: $-1.91$ at 100\% real, $-4.86$ at 50\%, $-6.08$ at 25\%, $-8.16$ at 10\%. With 93 real images against 3{,}068 synthetic ones the real set is outnumbered 33 to 1, and a frozen head cannot reconcile the two distributions --- it fits the synthetic majority. \emph{The identical pool that helps an unfrozen classifier harms a frozen one, and harms it most exactly where the paper claims augmentation should help most.}

\subsection{Supporting Analysis B: Generalization Curve}
To verify that our unfrozen classifier did not suffer from overfitting, we plotted train vs. test accuracy across 25 epochs.
\begin{itemize}[nosep, leftmargin=1.5em]
  \item \textbf{Result:} While training accuracy saturated near 100\% very early (by epoch 3--6), \textbf{test accuracy never degraded}. Instead, both the CNN-AD and CNN-SA test curves reached a stable plateau and remained there through epoch 25, with CNN-SA consistently resting above CNN-AD. Train saturation did not translate into test degradation, so early stopping was not required for a valid result here.
  \item \textbf{Protocol note:} the reported models use a \textbf{fixed 15 epochs chosen test-blind}. The per-epoch test curve above is a \emph{diagnostic} run only and was never used to select the epoch --- choosing an epoch by test score would be selecting on the test set and would bias the result. The apparent single-epoch peaks (CNN-AD 95.83\% at epoch 24, CNN-SA 97.40\% at epoch 6) are noise spikes at unrelated epochs, and are a good illustration of exactly why we did not tune on them. Stopping on a held-out \emph{validation} split would be a cleaner protocol still, and we note it as future work.
\end{itemize}


\section{Discussion} \label{sec:discussion}

\subsection{What Worked}
Stage 2 behaved as the Stage 1 diagnosis predicted. Unfreezing the encoder with discriminative learning rates adapted ImageNet features to CXR texture, raising the baseline on \emph{both} datasets, and it restored a real multi-seed augmentation lift (96.67\% on COVID, past the paper's 95\%). Track A also did its job on the generator: noise 1.0, DiffAugment, spectral norm and the projection discriminator more than halved FID and lifted the transfer probe from 65.5\% to 82.3\%, removing the brightness fingerprint.

\subsection{What Did Not Work}
\begin{itemize}[nosep, leftmargin=1.5em]
  \item \textbf{The +10-point magnitude did not reproduce.} We read this as available headroom rather than a broken method: a 90.62\% baseline has no room for a 10-point jump, and our scarcity sweep shows the lift growing as the baseline weakens.
  \item \textbf{Improving the GAN did not improve the classifier.} Under a frozen backbone even our best generator was flat or worse, and at a matched 300-epoch budget neither improved architecture separated from the baseline (Table~\ref{tab:generator_performance}). Every Track A intervention hit its stated target; none moved the number the paper's claim is stated in.
  \item \textbf{The two tracks do not compose.} We expected a better generator and a better classifier to add up. Capacity is a precondition for augmentation to help at all, but given capacity, generator quality had no reliable further effect --- and at the top of the range it turned negative.
\end{itemize}

\subsection{What We Learned}
Two knobs govern whether synthetic augmentation helps. \textbf{Capacity:} the downstream model needs trainable encoder capacity to absorb synthetic data; a frozen extractor can only shift a linear head, giving flat results or worse when swamped. \textbf{Headroom:} augmentation pays off in starved, low-baseline regimes and yields diminishing returns near ceiling. Beyond these, our central lesson is methodological: \textbf{generator-side metrics were \emph{inversely} related to downstream value across our four generators.} Judging a generator on its own is an unsafe proxy --- only a direct AD-vs-SA comparison tests an augmentation claim.

\subsection{Limitations and Future Work}
\begin{itemize}[nosep, leftmargin=1.5em]
  \item \textbf{Power.} CovidGAN results are 5-seed (3-seed frozen and scarcity), DDSM 2-seed. We report paired contrasts with standard deviations rather than formal tests, and mark contrasts that do not reach significance.
  \item \textbf{Anomaly A is unsettled.} A fine-scale dataset artifact would be indistinguishable from real pathology in our blurring test. Settling it needs an independent CXR dataset with both classes and no overlap with our sources --- harder than it sounds, since the IEEE collection is one of the paper's own COVID sources and several public benchmarks derive from it. This is the most valuable remaining check.
  \item \textbf{We could not recreate the paper's exact data.} No image-level manifest exists for their 403/721 selection, and the Kaggle database has been re-released since their March-2020 snapshot. Our deviation is one of \emph{version}, not source, and remains our leading explanation for the higher baseline.
  \item \textbf{Untested levers.} The two improved-GAN recipes are not yet unified across datasets, so the cross-dataset comparison is a parallel replication rather than one controlled experiment. The dilution hypothesis behind the reversal is unverified --- capping synthetic data at $1\times$ the real count, or augmenting only the minority class, would test it. Pretrain-then-finetune on the synthetic pool is an untried alternative to pooling. Finally, we used a fixed test-blind epoch count; validation-based early stopping would be cleaner.
\end{itemize}

\subsection{Reproducibility}
The GAN trained for the paper's full 2000 epochs ($\approx$18 hours on an Apple-silicon MPS GPU), checkpointing every 100 epochs; each detector trains in about a minute, which is what made the multi-seed grids here practical. All classifier runs and synthetic-pool generation are seeded and reproduce exactly. Every reported number is written to a \code{summary.json} recording seed, unfreeze depth and synthetic pool, so each table cell traces back to the run that produced it.

\begin{thebibliography}{10}
\bibitem{covidgan} A.~Waheed, M.~Goyal, D.~Gupta, A.~Khanna, F.~Al-Turjman, and P.~R.~Pinheiro, ``CovidGAN: Data Augmentation Using Auxiliary Classifier GAN for Improved Covid-19 Detection,'' \emph{IEEE Access}, vol.~8, pp.~91916--91923, 2020.
\bibitem{acgan} A.~Odena, C.~Olah, and J.~Shlens, ``Conditional Image Synthesis with Auxiliary Classifier GANs,'' in \emph{Proceedings of the International Conference on Machine Learning (ICML)}, 2017.
\bibitem{projection} T.~Miyato and M.~Koyama, ``cGANs with Projection Discriminator,'' in \emph{Proceedings of the International Conference on Learning Representations (ICLR)}, 2018.
\bibitem{vgg16} K.~Simonyan and A.~Zisserman, ``Very Deep Convolutional Networks for Large-Scale Image Recognition,'' in \emph{Proceedings of the International Conference on Learning Representations (ICLR)}, 2015.
\bibitem{fid} M.~Heusel, H.~Ramsauer, T.~Unterthiner, B.~Nessler, and S.~Hochreiter, ``GANs Trained by a Two Time-Scale Update Rule Converge to a Local Nash Equilibrium,'' in \emph{Advances in Neural Information Processing Systems (NeurIPS)}, 2017.
\bibitem{spectral} T.~Miyato, T.~Kataoka, M.~Koyama, and Y. Yoshida, ``Spectral Normalization for Generative Adversarial Networks,'' in \emph{Proceedings of the International Conference on Learning Representations (ICLR)}, 2018.
\bibitem{checkerboard} A.~Odena, V.~Dumoulin, and C.~Olah, ``Deconvolution and Checkerboard Artifacts,'' \emph{Distill}, 2016.
\bibitem{diffaugment} S.~Zhao, Z.~Liu, J.~Lin, J.-Y. Zhu, and S. Han, ``Differentiable Augmentation for Data-Efficient GAN Training,'' in \emph{Advances in Neural Information Processing Systems (NeurIPS)}, 2020.
\bibitem{kaggle_covid} T.~Rahman, M.~E.~H.~Chowdhury et al., ``COVID-19 Radiography Database,'' Kaggle.
\bibitem{ddsm} R.~S.~Lee, F.~Gimenez, A.~Hsu et al., ``A curated mammography data set for use in computer-aided detection and diagnosis research'' (CBIS-DDSM), \emph{Scientific Data}, 2017.
\end{thebibliography}

\end{document}
